% ---------- Datei: chapters/07_projektabschluss.tex ----------
\chapter{Projektabschluss}

\section{Soll-Ist-Vergleich}
Das Projekt hat alle im Antrag definierten Ziele erfolgreich erfüllt. Eine zentral verwaltete Hub-and-Spoke-Architektur mit dem Google Cloud Network Connectivity Center (NCC) hat die unübersichtliche VPN-Struktur abgelöst. Das neue Setup mit HA VPN und dynamischem BGP-Routing läuft stabil und bietet die geforderte Ausfallsicherheit. Auch die Umsetzung der Sicherheitsvorgaben (Default-Deny) und des zentralen Monitorings hat wie geplant stattgefunden.

\begin{table}[!ht]
\centering
\begin{tabularx}{\textwidth}{|X|X|X|}
\hline
\textbf{Ziel} & \textbf{Soll} & \textbf{Ist} \\ \hline
Zentrale Steuerung & NCC-Hub & Erreicht \\ \hline
Standortanbindung & HA-VPN & Erreicht \\ \hline
Routing & BGP (dynamisch) & Erreicht \\ \hline
Sicherheit & Firewall Default-Deny & Erreicht \\ \hline
Monitoring & Zentrale Überwachung & Erreicht \\ \hline
\end{tabularx}
\caption{Soll-Ist-Vergleich der Projektziele}
\end{table}

\section{Zeitaufwand}
Die tatsächliche Projektzeit hat die geplanten 40 Stunden mit 39,5 Stunden leicht unterschritten. Dabei hat es kleinere Verschiebungen gegeben: Die Vorbereitungsphase hat etwas mehr Zeit beansprucht (+1,5\,h), um das Skripting für das Deployment detailliert auszuarbeiten. Diese Vorarbeit hat sich jedoch ausgezahlt, da die anschließende Cloud-Bereitstellung dadurch deutlich schneller und nahezu fehlerfrei abgelaufen ist (-2,0\,h).

\begin{table}[!ht]
\centering
\begin{tabular}{|p{0.34\textwidth}|c|c|}
\hline
\textbf{Projektphase} & \textbf{Geplant (h)} & \textbf{Ist (h)} \\ \hline
Planung & 8 & 8 \\ \hline
Vorbereitung & 6 & 7,5 \\ \hline
Durchführung & 12 & 10 \\ \hline
Test und QA & 5 & 5 \\ \hline
Übergabe & 3 & 3 \\ \hline
Dokumentation & 6 & 6 \\ \hline
\textbf{Gesamt} & \textbf{40} & \textbf{39,5} \\ \hline
\end{tabular}
\caption{Gegenüberstellung des geplanten und tatsächlichen Zeitaufwands}
\end{table}

\section{Fazit und Ausblick}
Die Umstellung auf eine zentrale NCC-Architektur in der Google Cloud reduziert den operativen Aufwand im Netzwerk deutlich. Durch das automatisierte BGP-Routing skaliert die Infrastruktur optimal. Neue Standorte lassen sich so in Zukunft schnell und nach einem festen Standard anbinden.

\subsection*{Einfluss und Einsatz von KI}
Die Nutzung von Künstlicher Intelligenz (\textbf{Google Gemini 3.1 Pro}) hat das Projekt als effizientes Hilfsmittel begleitet. Besonders bei der Validierung neuer \texttt{gcloud}-Syntax für NCC-Spokes und bei der Strukturierung des \LaTeX-Codes für diese Dokumentation hat das Tool den zeitlichen Rechercheaufwand reduziert.

Trotz der KI-Unterstützung haben alle architektonischen Entscheidungen, die fundierte Fehlersuche (Troubleshooting) sowie die praktische Validierung eine eigenständige fachliche Leistung erfordert. Das Projekt hat verdeutlicht, dass KI ein wertvolles Werkzeug zur Effizienzsteigerung ist, belastbares IT-Fachwissen für die Bewertung der Ergebnisse jedoch zwingend voraussetzt.